% This file compiles with both LuaLaTeX and XeLaTeX
\documentclass[11pt]{article}

% Change "review" to "final" to generate the final (sometimes called camera-ready) version.
% Change to "preprint" to generate a non-anonymous version with page numbers.
\usepackage[preprint]{acl}

% This is not strictly necessary, and may be commented out,
% but it will improve the layout of the manuscript,
% and will typically save some space.
 \usepackage{microtype}

% If the title and author information does not fit in the area allocated, uncomment the following
%
%\setlength\titlebox{<dim>}
%
% and set <dim> to something 5cm or larger.

% These font selection commands work with
% LuaLaTeX and XeLaTeX, but not pdfLaTeX.
\usepackage[english,bidi=default]{babel} % English as the main language.
\babelfont{rm}{TeXGyreTermesX} % similar to Times
%%% include whatever languages you need below this line
\babelprovide[import]{vietnamese}
\babelfont[*vietnamese]{rm}{TeXGyreTermesX}

\usepackage{graphicx}
\usepackage{booktabs}
\usepackage{multirow}
\usepackage{amsmath}
\usepackage{amsfonts}
\usepackage{bm}


\title{Multi-Aspect Heterogeneous Graph Neural Networks for \\ Vietnamese News Recommendation}

\author{Student Name \\
  University Name \\
  \texttt{email@domain} \\\And
  Advisor Name \\
  University Name \\
  \texttt{email@domain} \\}

\begin{document}

\maketitle

\begin{abstract}
    News recommendation in Vietnamese presents unique challenges due to linguistic complexity, sparse user interactions, and the "cold-start" problem. In this work, we propose a comprehensive breakdown of the recommendation pipeline, leveraging Multi-Aspect Heterogeneous Graph Neural Networks (MA-HGN) and Contrastive Learning (CL). We introduce a novel heterogeneous graph construction (G4) incorporating users, articles, categories, and authors, enriched with weighted edges for engagement and recency. We benchmark state-of-the-art Collaborative Filtering (CF) models, including NGCF, LightGCL, SimGCL, and XSimGCL, alongside our custom implementations of MA-HCL and SimMAHGN. Our extensive experiments on a real-world VnExpress dataset utilizing 10k users and 3.7k articles demonstrate that our hybrid approach, combining SimGCL with VN-SBERT embeddings, achieves superior performance (NDCG@10 of 0.1895) compared to traditional methods. We also present a detailed ablation study on edge types and embedding models, highlighting the critical role of semantic enrichment in news recommendation.
\end{abstract}

\section{Introduction}

In the era of information overload, personalized news recommendation systems (NRS) are essential for helping users discover relevant content. While significant progress has been made in English-language NRS, Vietnamese NRS remains under-explored, primarily due to the lack of large-scale datasets and the complexity of the language.

Traditional approaches often rely on collaborative filtering (CF), which suffers from data sparsity, or content-based (CB) methods, which may lack serendipity. Graph Neural Networks (GNNs) have emerged as a powerful tool to bridge this gap, capable of modeling complex interactions between users and items. However, most existing GNN-based NRS model the data as a bipartite user-item graph, ignoring rich auxiliary information such as article categories, authors, and temporal dynamics.

In this paper, we address these limitations by modeling the news ecosystem as a \textbf{Heterogeneous Information Network (HIN)}. Our key contributions are:
\begin{enumerate}
    \item \textbf{Heterogeneous Graph Construction:} We construct extensive graph variants (G1-G4), integrating Author and Category nodes, and introducing weighted edges based on user engagement (comments, replies) and article freshness.
    \item \textbf{Advanced Modeling:} We implement and adapt \textbf{SimMAHGN} (Simple Multi-Aspect Heterogeneous Graph Network) and \textbf{MA-HCL} (Multi-Aspect Heterogeneous Contrastive Learning) to the Vietnamese context.
    \item \textbf{Comprehensive Benchmarking:} We provide a rigorous evaluation of various embeddings (VN-SBERT, PhoBERT, BGE-M3) and CF models (SimGCL, XSimGCL, LightGCL).
    \item \textbf{Deployment:} We develop an interactive Streamlit application demonstrating real-time hybrid recommendation with A/B testing capabilities.
\end{enumerate}

\section{Related Work}

\subsection{Graph Neural Networks for Recommendation}
Early graph-based recommendation methods like NGCF \citet{wang2019neural} explicitly modeled high-order connectivities in the user-item graph. However, the heavy computational cost of non-linear activations led to the proposal of LightGCL \citet{he2020lightgcn}, which simplified the GCN architecture by removing linear transformations and non-linearities. Recently, SimGCL \citet{yu2022graph} integrated contrastive learning (CL) into LightGCL, using random noise perturbations to enforce representation consistency, significantly improving robustness against sparse data.

\subsection{Heterogeneous Graph Recommendation}
Heterogeneous GNNs (HGNNs) extend GNNs to graphs with multiple node and edge types. Methods like HAN and HGT utilize meta-paths to aggregate neighbors. In news recommendation, \citet{liu2021heterogeneous} proposed modeling news as a heterogeneous graph including entities and categories, which inspired our MA-HCL approach. Unlike previous works that focus solely on English datasets (MIND), we adapt these architectures for Vietnamese news by integrating language-specific embeddings.

\subsection{Pretrained Language Models in SRS}
PLMs like BERT \citet{devlin2018bert} have revolutionized content understanding. In the Vietnamese context, PhoBERT and VN-SBERT \citet{reimers2019sentence} provide strong baselines. Our work systematically evaluates these embeddings within a hybrid NRS framework.

\section{Methodology}

\subsection{Data Collection and Preprocessing}
We crawled data from VnExpress, a major Vietnamese news portal, resulting in 3,685 articles and 42,809 user interactions. The data includes:
\begin{itemize}
    \item \textbf{Users:} Anonymized user IDs and interaction history.
    \item \textbf{Articles:} Title, short description, content, publication date, and author.
    \item \textbf{Interactions:} User comments and replies, serving as a proxy for engagement.
\end{itemize}

\subsection{Heterogeneous Graph Construction}
We define a heterogeneous graph $G = (V, E)$ where $V = \mathcal{U} \cup \mathcal{I} \cup \mathcal{C} \cup \mathcal{A}$ consists of user, article, category, and author nodes.
We experimented with four graph variants (G1-G4) to isolate the impact of different signals. The most comprehensive graph, \textbf{G4}, includes:
\begin{enumerate}
    \item \textbf{User-Article Edges ($E_{ui}$):} We use a composite weighting scheme that captures both engagement quality and temporal relevance:
          \begin{equation}
              w_{ui} = (1 + \log(1 + r)) \times e^{-\lambda \Delta t}
          \end{equation}
          where $r$ is the reaction count on the user's comment, $\Delta t$ is the number of days since the interaction, and $\lambda = 0.01$ is the daily decay rate. The logarithmic term dampens the influence of viral outliers while the exponential decay prioritizes recent interactions.
    \item \textbf{Article-Category Edges ($E_{ic}$):} connecting articles to their topics (e.g., ``Sports'', ``Business'').
    \item \textbf{Article-Author Edges ($E_{ia}$):} connecting articles to their writers.
    \item \textbf{User-User Edges ($E_{uu}$):} Implicit social links based on direct replies in the comment section.
\end{enumerate}

\subsection{Model Architecture}

\subsubsection{SimGCL \& XSimGCL}
We adopted SimGCL \citet{yu2022graph} as a strong baseline foundation. SimGCL simplifies Graph Convolutional Networks (GCN) by removing non-linearities and introducing contrastive learning (CL) noise to regularize embedding representations.

Let $E^{(k)}$ be the embedding matrix at layer $k$. The propagation rule is:
\begin{equation}
    E^{(k)} = \frac{1}{2} (D^{-\frac{1}{2}} A D^{-\frac{1}{2}}) E^{(k-1)}
\end{equation}
SimGCL augments the representation learning with a contrastive auxiliary task:
\begin{equation}
    \mathcal{L}_{CL} = \sum_{i \in \mathcal{U} \cup \mathcal{I}} - \log \frac{\exp(sim(e_i^{'}, e_i^{''}) / \tau)}{\sum_{j \neq i} \exp(sim(e_i^{'}, e_j^{''}) / \tau)}
\end{equation}
where $e_i^{'}$ and $e_i^{''}$ are perturbed embeddings with added random noise.

\subsubsection{MA-HCL (Multi-Aspect Heterogenous CL)}
Adapted from \citet{liu2021heterogeneous}, MA-HCL handles heterogeneity using meta-path based aggregation. For each meta-path $\Phi$ (e.g., User-Article-Author-Article), we compute a distinct embedding $h_u^{\Phi}$ for user $u$.
We employ a \textbf{hierarchical attention mechanism} to fuse embeddings from different meta-paths:
\begin{equation}
    h_u = \sum_{\Phi} \alpha_{\Phi} h_u^{\Phi}
\end{equation}
where $\alpha_{\Phi}$ is the learned importance of meta-path $\Phi$. To improve robustness, we apply contrastive learning on the meta-path views, encouraging consistency between the global representation and aspect-specific representations.

\subsubsection{Hybrid Inference Strategy}
To address the cold-start problem where graph connectivity is sparse or non-existent (new users/items), we implement a hybrid strategy.
Let $s_{CF}(u, i)$ be the score from the CF model (SimGCL/MA-HCL).
Let $s_{CB}(u, i)$ be the cosine similarity between user history and candidate item embedding (VN-SBERT).
The final score is:
\begin{equation}
    s_{final} = \beta \cdot \sigma(s_{CF}) + (1-\beta) \cdot s_{CB}
\end{equation}
where $\beta$ is a dynamic weight parameter regulated by user activity level.

\section{Experiments}

\subsection{Experimental Setup}
\textbf{Dataset Split:} We use a chrono-holdout split, training on the first 80\% of interactions and testing on the subsequent 20\% to mimic real-world scenarios.

\textbf{Metrics:} We report Recall@K and NDCG@K ($K=10, 20$).

\textbf{Implementation:} Models are implemented in PyTorch Geometric. We use Adam optimizer with learning rate 0.001, batch size 2048, and train for 50 epochs.

\subsection{Main Results}
Table \ref{tab:results} summarizes the performance of our models.

\begin{table}[h]
    \centering
    \small
    \begin{tabular}{lcc}
        \hline
        \textbf{Model}             & \textbf{Recall@10} & \textbf{NDCG@10} \\
        \hline
        \multicolumn{3}{l}{\textit{Baselines}}                             \\
        TF-IDF                     & 0.0465             & 0.0250           \\
        LightGCL                   & 0.0482             & 0.0300           \\
        SimGCL                     & 0.0645             & 0.0412           \\
        XSimGCL                    & 0.0782             & 0.0452           \\
        \hline
        \multicolumn{3}{l}{\textit{Ours (Heterogeneous)}}                  \\
        SimMAHGN                   & 0.1240             & 0.0769           \\
        \textbf{MA-HCL}            & \textbf{0.1502}    & \textbf{0.1024}  \\
        \hline
        \multicolumn{3}{l}{\textit{Hybrid (CF + VN-SBERT)}}                \\
        \textbf{SimGCL + VN-SBERT} & \textbf{0.2415}    & \textbf{0.1895}  \\
        \hline
    \end{tabular}
    \caption{Performance comparison on VnExpress dataset.}
    \label{tab:results}
\end{table}

\textbf{SimGCL + VN-SBERT} achieved the highest performance (R@10=0.2415), confirming that semantic enrichment is critical for news recommendation where item text carries significant signal. Among pure collaborative filtering methods, our \textbf{MA-HCL} model outperformed the strongest baseline (XSimGCL) by a large margin (0.1502 vs 0.0782), validating the effectiveness of our heterogeneous graph construction and multi-aspect contrastive learning.

\subsection{Ablation Studies}

\subsubsection{Impact of Text Embeddings}
We compared various text embeddings for the content-based component.
\begin{table}[h]
    \centering
    \small
    \begin{tabular}{lc}
        \hline
        \textbf{Embedding Model} & \textbf{Relative Perf.} \\
        \hline
        TF-IDF                   & Baseline (1.0x)         \\
        PhoBERT                  & 1.15x                   \\
        \textbf{VN-SBERT}        & \textbf{1.42x}          \\
        BGE-M3                   & 1.35x                   \\
        \hline
    \end{tabular}
    \caption{Impact of embedding choice on CB Performance.}
    \label{tab:embeddings}
\end{table}
VN-SBERT proved to be the most effective (1.42x improvement over TF-IDF), likely due to its sentence-level optimization which aligns well with news titles and descriptions compared to the token-level masking task of PhoBERT.

\subsubsection{Impact of Graph Structure}
We analyzed the contribution of incremental graph complexity:
\begin{itemize}
    \item \textbf{G1 (User-Article):} Simple bipartite baseline.
    \item \textbf{G2 (+ Social):} Weighted edges from comments and replies. +248\% gain over G1.
    \item \textbf{G3 (+ Category Hubs):} Semantic shortcuts via topic nodes. +984\% total gain.
    \item \textbf{G4 (Full Heterogeneous):} Adding Author nodes and multi-aspect fusion for maximum schematic richness.
\end{itemize}

\section{Conclusion}
We presented a comprehensive study of Graph Neural Networks for Vietnamese news recommendation. Our work underscores the value of heterogeneous graphs and contrastive learning. The proposed hybrid framework, combining efficient GNNs with high-quality semantic embeddings, offers a robust solution for personalized news delivery. Future work will explore dynamic graph updates and large language model (LLM) integration for explanation generation.

\section*{Limitations}
Our dataset, while representative, is smaller than standard academic benchmarks (MIND, Adressa). The "cold-start" strategy relies heavily on content similarity, which may not capture changing user interests immediately.

\bibliography{custom}

\end{document}
